% Page budget: 1.55 pages | Owns: negation, OBC derivation, additional-state scope, implementation, and true-parameter condition; model M_OBC.
% WRITING GUIDE
% Purpose: Explain negation, derive the final state-level orthogonal-by-construction map, and delimit its interpretability guarantee.
% Sources: README "Source map per subsection", rows IV-A, IV-B and IV-C; mandatory papers and the reference gate as in the README.
% Research-plan anchor: Preserve Aspect 3's goal of protecting physical interpretability; state clearly how the final construction differs from the proposed regularisation.
% Current decisions: Reduced identifiable coordinates, data-derived reference tuples, a tangent-only one-step physical-state projection, per-epoch refresh, and no affine comparison arm.
% Choices to justify: Protected parameter coordinates, reference-set construction, projected rows, rank tolerance, exclusion of the affine offset, refresh frequency, and treatment of additional states.
% Cautions: docs/orthogonal-projection-plan.md describes an older soft-penalty route; one-step state orthogonality is not trajectory-output orthogonality or guaranteed true recovery.
% Planned figures: New geometric projection figure with a physical and additional-state partition inset; show the remaining uniqueness limitation as clearly as the protected direction.
%   (Figure plan, not yet drawn: decomposition of F_R into its part in range(Phi_R) and its orthogonal part at
%   the expansion point, a curved surface marking that the tangent is local; inset with the corrected physical
%   rows, the uncorrected added-state rows and their route back through f_aug.)
% Section is finished when: The protected space, coordinate frame, gradient treatment, added-state scope, and true-recovery condition are explicit.
%
% MUST ESTABLISH (inferred by the sandbox writer, cycle 02, 2026-10-09, from the header, the README ownership
%   row, Section III's handover and Documentation/THESIS-RESULTS.md step 5; NOT approved by Dirk)
% Contribution delivered: Introduction bullet 3, state-level OBC for the identifiable tangent space and its
%   true-parameter recovery condition. Adopted: Gyorok 2026 coefficient, projection, Lemma 3, prediction rule.
%   Adapted: regressor replaced by the RK4 transition sensitivity at an expansion point refreshed per epoch
%   (linearisation as in Gyorok 2025 Sec. 4); reference states reconstructed from positions instead of a
%   baseline simulation. This work: the state-level LPV MIMO realisation, the local uniqueness statement, the
%   state-level recovery condition, the scope.
% IV-A 1. Negation (Kessels' term) made precise to first order; Section III already established
%         non-uniqueness, do not repeat it. 2. Gantry directions are generalized forces with the structure
%         of M(Y), C, K through M(Y)^-1.
% IV-B 3. Construction on data-derived tuples (x_a = 0), coefficient once per objective with gradient
%         through the solve, basis per epoch, offset not protected, added-state rows uncorrected; model
%         M_OBC = M_PS2 with the corrected transition, same criterion V. Differs from the research plan's
%         penalty (no weight).
% IV-C 4. Scope: stacked, one-step, local, x_a = 0, normalised metric; Gyorok's Thms 16 and 18 do not
%         transfer. 5. Local uniqueness (short statement and proof, README Math standard item 5).
%         6. Unique is not true: error Phi^+ delta*, recovery iff Phi' delta* = 0; simulation-only check.
% Left out on purpose: the affine [J | c] arm (not in the campaign, THESIS-RESULTS Sec. 6); the parameter
%   prior (Section III MUST ESTABLISH: no parameter regularisation anywhere; Setup says it is off); the
%   hand-written forward sensitivity and compile seam (implementation, same model either way); the
%   training-coordinate Jacobian (diffeomorphic reparametrisation, same corrected model).
% Moved out earlier and kept out: measurement details of the pointwise no-go result and the Coulomb
%   inner-product argument (Section VII); rank, condition number and column correlations of the stack
%   (Results); rho* = 0.205 and the 4.96 % floor (earlier absorber-only dataset, not the thesis truth).
\section{Interpretability by Orthogonal Construction}\label{sec:obc}
\ifdraft
\begin{itemize}
\item Given the identified model of Section~III, the goal is a unique estimate of $\theta_{\mathrm{base}}$; the section names its three parts (style profile)
\item Contribution share: OBC carried from linear-in-parameter input-output models to the state transition of the self-scheduled gantry baseline, with its uniqueness and recovery conditions (Introduction, third contribution; Gy\"or\"ok 2026 Sec.~6)
\end{itemize}
\fi

Given the augmented model \eqref{eq:aug_transition}, identified by
\eqref{eq:aug_loss}, our goal is a unique estimate of
$\theta_{\mathrm{base}}$, so that it keeps its physical meaning
(Section~\ref{sec:aug_structure}).
To this end, Section~\ref{sec:obc_negation} states how the network makes the
estimate non-unique, Section~\ref{sec:obc_map} removes that mechanism by an
orthogonal-by-construction (OBC) correction of the network, and
Section~\ref{sec:obc_scope} states what the correction guarantees and when the
estimate equals the true parameters.
Gy\"or\"ok et al.\ define OBC for input-output models that are linear in their
parameters and name state-space baselines without full-state measurement as
open~\cite[Sec.~6]{gyorok2026obc}.
We extend it to the state transition of the self-scheduled gantry baseline.

\subsection{Negation of the Physical Parameters}\label{sec:obc_negation}
\ifdraft
\begin{itemize}
\item The baseline alone is not ambiguous (Section~II-C); the network makes the split non-unique (Section~III-A); Kessels calls the effect negation (Kessels Sec.~5.3.1.3)
\item Shown for linear baselines by Gy\"or\"ok (2025 Ex.~1, 2026 Ex.~2); for the nonlinear transition it has an exact first-order form with the sensitivity $\Phi$ (this work, linearisation as in Gy\"or\"ok 2025 Eq.~(15))
\item On the gantry the negated directions are generalized forces with the structure of $M(Y)$, $C$, $K$, scaled by $M(Y)^{-1}$ and so dependent on $Y$ (this work)
\end{itemize}
\fi

The baseline alone is not ambiguous, because its ten combinations are
structurally identifiable (Section~\ref{sec:params}).
The network makes the split non-unique (Section~\ref{sec:aug_structure}).
Kessels calls this effect, learned terms that partially cancel the
first-principles terms, negation~\cite[Sec.~5.3.1.3]{kessels2025ai}.
Gy\"or\"ok et al.\ show it for a baseline linear in its parameters with a
linear network layer~\cite[Ex.~1]{gyorok2025l4dc},~\cite[Ex.~2]{gyorok2026obc}.
For the nonlinear gantry transition, negation has an exact first-order form.
For an expansion point $\bar\theta_{\mathrm{base}}$ and a displacement
$p\in\R^{10}$,
\begin{equation}
\begin{aligned}
 &f_{\mathrm{base}}(\bar\theta_{\mathrm{base}}+p,\tilde x,u)
 +\bigl[f_{\mathrm{aug}}(\theta_{\mathrm{aug}},\hat x,u)-\Phi(\tilde x,u)\,p\bigr]\\
 &\quad=f_{\mathrm{base}}(\bar\theta_{\mathrm{base}},\tilde x,u)
 +f_{\mathrm{aug}}(\theta_{\mathrm{aug}},\hat x,u)+O(\norm{p}^2),
\end{aligned}
\label{eq:obc_negation}
\end{equation}
where
$\Phi(\tilde x,u)=\partial f_{\mathrm{base}}(\theta_{\mathrm{base}},\tilde x,u)/\partial\theta_{\mathrm{base}}
|_{\bar\theta_{\mathrm{base}}}$, with $\Phi:\R^6\times\R^3\to\R^{6\times10}$,
is the sensitivity of the one-step transition to the ten combinations.
A network that adds $-\Phi p$ therefore hides any displacement $p$ of
$\theta_{\mathrm{base}}$ from the data.

On the gantry, the columns of $\Phi$ are generalized forces passed through
$M(Y)^{-1}$: the mass combinations act through $\ddot q$, the damping
combinations through $\dot q$, and $k_{b,\Sigma}$ through $q$.
Differentiating \eqref{eq:eom} at fixed $(q,\dot q,u)$ gives
\begin{equation}
 \frac{\partial\ddot q}{\partial\theta_{\mathrm{base},j}}
 =-M(Y)^{-1}\Bigl[\frac{\partial M(Y)}{\partial\theta_{\mathrm{base},j}}\,\ddot q
 +\frac{\partial C}{\partial\theta_{\mathrm{base},j}}\,\dot q
 +\frac{\partial K}{\partial\theta_{\mathrm{base},j}}\,q\Bigr],
\label{eq:obc_gantry_sens}
\end{equation}
where $\theta_{\mathrm{base},j}$ is the $j$-th entry of
$\theta_{\mathrm{base}}$, and the RK4 stages carry this acceleration
sensitivity into $\Phi$.
A negating network term is therefore an added inertia, damping or stiffness
force whose direction depends on the payload position.

\subsection{Orthogonal-by-Construction for the Gantry}\label{sec:obc_map}
\ifdraft
\begin{itemize}
\item OBC fits the network output to the baseline range on a data set and subtracts the fit; it replaces the projection penalty and needs no weight (adopted, Gy\"or\"ok 2026 Eqs.~(8), (10), Lemma~3; Gy\"or\"ok 2025 Eq.~(13))
\item The gantry transition has no regressor, so $\Phi$ at $\bar\theta_{\mathrm{base}}$ replaces it; no state is measured and a baseline simulation drifts on $X$ and $Y$, so the reference states are reconstructed from the positions, on $\bar x=0$ (adapted from Gy\"or\"ok 2025 Sec.~3, Eq.~(15); D-192)
\item Stacks over all six physical rows, coefficient and stacked orthogonality on the reference set (Gy\"or\"ok 2026 Lemma~3; D-192)
\item Full column rank required; SVD solve, rank tolerance and rank-loss rule (Gy\"or\"ok 2026 Assumption~1; D-192, code)
\item Corrected model: correction on the physical rows, $g_{\mathrm{aug}}$ uncorrected (Gy\"or\"ok 2026 Eq.~(13); this work)
\item $\mathcal M_{\mathrm{OBC}}$: PS2 start, criterion $V$, gradient through the solve (this work; D-192, D-240)
\item Coefficient once per evaluation of $V$, basis rebuilt per epoch at the current estimate; validation and the reported model synchronised (adapted from Gy\"or\"ok 2025 Sec.~4; D-192, D-198)
\item The offset of the expansion is not protected (D-192; unlike Gy\"or\"ok 2025 Eq.~(18))
\end{itemize}
\fi

Gy\"or\"ok et al.\ remove negation by construction.
For a baseline linear in its parameters, they fit the stacked network output to
the range of the stacked regressor on a data set and subtract the
fit~\cite[Eqs.~(8), (10)]{gyorok2026obc}.
The corrected network is then orthogonal to every baseline direction on that
set~\cite[Lemma~3]{gyorok2026obc}.
This replaces the projection penalty of their earlier
method~\cite[Eq.~(13)]{gyorok2025l4dc}, which needs a weight that trades
orthogonality against accuracy.

Two properties of the gantry prevent a direct transfer.
The transition is rational in $\theta_{\mathrm{base}}$ through $M(Y)^{-1}$ and
has no regressor.
Following the linearisation in~\cite[Eq.~(15)]{gyorok2025l4dc}, we use
$\Phi$ in its place.
The stacks also need states, while only positions are measured.
Gy\"or\"ok et al.\ obtain states by simulating the baseline on the training
data~\cite[Sec.~3]{gyorok2025l4dc}.
On the free $X$ and $Y$ axes such a simulation drifts
(Section~\ref{sec:aug_closed_loop}).
We therefore reconstruct one reference state from every training sample $k$
whose stencil lies inside its record.
The damping directions multiply the velocity, so its estimate uses a
fourth-order central difference.
The recorded input already contains the feedback, so the controller does not
enter.
In the coordinates of \eqref{eq:aug_norm},
\begin{equation}
\begin{aligned}
 q^{\mathrm r}_k&=P^{-\top}y_k,\qquad
 \dot q^{\mathrm r}_k=\frac{-q^{\mathrm r}_{k+2}+8q^{\mathrm r}_{k+1}
 -8q^{\mathrm r}_{k-1}+q^{\mathrm r}_{k-2}}{12\,T_s},\\
 \tilde x^{\mathrm r}_k&=T_x\bigl(\col(q^{\mathrm r}_k,\dot q^{\mathrm r}_k)-\mu_x\bigr),
 \qquad k\in\mathcal K,
\end{aligned}
\label{eq:obc_tuples}
\end{equation}
where $q^{\mathrm r}_k,\dot q^{\mathrm r}_k\in\R^3$ are the reconstructed
positions and velocities, $\tilde x^{\mathrm r}_k\in\R^6$ the reference state,
$\mathcal K$ the set of these samples, and
$\mathcal R=\{(\tilde x^{\mathrm r}_k,u_k)\}_{k\in\mathcal K}$ the reference
set of $N_{\mathcal R}=|\mathcal K|$ tuples with the recorded, normalised
input $u_k$.
Each tuple evaluates $M(Y)$ at its own measured $Y$, so $\mathcal R$ covers
the scheduling range of the training records.
\todo{Missing: the reference velocities in the noisy thesis mode. \eqref{eq:obc_tuples} differentiates the anti-alias filtered noisy positions, and the guard in \texttt{GD/obc\_gantry.py::build\_reference\_set} checks only \texttt{cfg.snr}, which the thesis data leave unset (sources.md trap~7). Candidate: state the velocity noise of the noisy versus the noise-free reference set in Section~\ref{sec:setup}, or build the tuples from the noise-free twin and say so (D-225, D-232).}

The network writes all six physical rows (Section~\ref{sec:aug_structure}),
so the stacks hold all six.
With the added states set to zero, the construction on $\mathcal R$ reads
\begin{equation}
\begin{aligned}
 \Phi_{\mathcal R}&=\col\bigl(\Phi(\tilde x^{\mathrm r}_k,u_k)\bigr)_{k\in\mathcal K}
 \in\R^{6N_{\mathcal R}\times10},\\
 F_{\mathcal R}(\theta_{\mathrm{aug}})&=\col\bigl(f_{\mathrm{aug}}(\theta_{\mathrm{aug}},
 \col(\tilde x^{\mathrm r}_k,0),u_k)\bigr)_{k\in\mathcal K},\\
 \theta_{\mathrm{aux}}&=\Phi_{\mathcal R}^{+}F_{\mathcal R}(\theta_{\mathrm{aug}}),\\
 0&=\Phi_{\mathcal R}^{\top}\bigl(F_{\mathcal R}(\theta_{\mathrm{aug}})
 -\Phi_{\mathcal R}\theta_{\mathrm{aux}}\bigr),
\end{aligned}
\label{eq:obc_coef}
\end{equation}
where $F_{\mathcal R}\in\R^{6N_{\mathcal R}}$ is the stacked network output,
$(\cdot)^{+}$ the Moore-Penrose pseudoinverse, and
$\theta_{\mathrm{aux}}\in\R^{10}$ the auxiliary coefficient.
By~\cite[Lemma~3]{gyorok2026obc}, the last line holds for every
$\theta_{\mathrm{aug}}$ while $\Phi_{\mathcal R}$ is fixed, so no network
update moves the corrected stack into $\ran\Phi_{\mathcal R}$.

The construction requires full column rank of $\Phi_{\mathcal R}$, the
counterpart of~\cite[Assumption~1]{gyorok2026obc}: the reference set must
separate the ten directions.
$\theta_{\mathrm{aux}}$ is computed from an economy SVD of
$\Phi_{\mathcal R}$, which avoids squaring its condition number.
The rank is decided with the tolerance
$\max(6N_{\mathcal R},10)\,\epsilon\,\norm{\Phi_{\mathcal R}}_2$, with
$\epsilon$ the machine precision.
A rebuild that loses rank keeps the previous basis.
A second consecutive loss stops the training.

The corrected network enters the state update as in the prediction rule
of~\cite[Eq.~(13)]{gyorok2026obc},
\begin{equation}
\begin{aligned}
 \tilde x_{k+1}&=f_{\mathrm{base}}(\theta_{\mathrm{base}},\tilde x_k,u_k)
 +f_{\mathrm{aug}}(\theta_{\mathrm{aug}},\hat x_k,u_k)
 -\Phi(\tilde x_k,u_k)\,\theta_{\mathrm{aux}},\\
 \bar x_{k+1}&=g_{\mathrm{aug}}(\theta_{\mathrm{aug}},\hat x_k,u_k),\qquad
 \hat y_k=h_{\mathrm{base}}(\tilde x_k),
\end{aligned}
\label{eq:obc_model}
\end{equation}
where $\Phi$ is evaluated at $\bar\theta_{\mathrm{base}}$ and the current
state, while $f_{\mathrm{base}}$ uses the current $\theta_{\mathrm{base}}$.
The baseline has no rows for $\bar x$, so $g_{\mathrm{aug}}$ has no baseline
direction to be orthogonal to and stays uncorrected.

We call the model \eqref{eq:obc_model} with encoder \eqref{eq:aug_encoder},
initialised by PS2 (Section~\ref{sec:aug_ps2}) and identified by
\eqref{eq:aug_loss} with $u=\hat u$, $\mathcal M_{\mathrm{OBC}}$.
It is trained over the same $\vartheta$ as $\mathcal M_{\mathrm{PS2}}$.
$\theta_{\mathrm{aux}}$ is not a trained parameter but a function of
$\theta_{\mathrm{aug}}$.
The gradient of $V$ therefore passes through the solve \eqref{eq:obc_coef}.

The coefficient and the basis are refreshed at different rates.
$\theta_{\mathrm{aux}}$ is solved once per evaluation of $V$ from the current
network, since a coefficient kept longer is orthogonal only for the network it
was solved from.
The expansion point has to follow the estimate.
Gy\"or\"ok et al.\ update it at every evaluation of the cost or fix it at the
nominal parameters~\cite[Sec.~4]{gyorok2025l4dc}.
Updating rebuilds $\Phi_{\mathcal R}$ at every evaluation.
A fixed point protects the directions at a parameter value that joint
estimation leaves.
We rebuild $\Phi_{\mathcal R}$ at each epoch boundary at the current
$\theta_{\mathrm{base}}$ and hold it, without gradient, for the epoch.
Validation uses the coefficient of the current network.
The reported model has its basis rebuilt once at its selected
$\hat\theta_{\mathrm{base}}$.

The protected space excludes the offset of the expansion.
The nonlinear method of~\cite[Eq.~(18)]{gyorok2025l4dc} also protects
$f_{\mathrm{base}}(\bar\theta_{\mathrm{base}},\tilde x,u)
-\Phi(\tilde x,u)\bar\theta_{\mathrm{base}}$, whereas we follow~\cite[Eq.~(8)]{gyorok2026obc},
which protects the regressor range only.
The offset multiplies no estimated parameter, so it carries no parameter
ambiguity.
A one-step transition is dominated by the propagation of the state itself, so
protecting the offset would also remove network directions that negate no
parameter.
\todo{Missing: supervisor confirmation of the protected set. D-190 records the supervisors' condition as non-overlap with the baseline output, which the offset would protect; D-192 protects $\ran\Phi_{\mathcal R}$ only. Candidate: keep $\ran\Phi_{\mathcal R}$ (parameter uniqueness,~\cite[Eq.~(8)]{gyorok2026obc}) and state the offset overlap as unprotected (Section~\ref{sec:obc_scope}).}

\subsection{Scope and Condition for True Parameter Recovery}\label{sec:obc_scope}
\ifdraft
\begin{itemize}
\item The orthogonality is stacked over $\mathcal R$, one step, local at $\bar\theta_{\mathrm{base}}$, on $\bar x=0$, in the normalised metric (D-192, D-202)
\item Not protected: the part of $f_{\mathrm{aug}}$ acting through nonzero $\bar x$ (active after PS2), multistep and closed-loop responses; Gy\"or\"ok's Thms.~16, 18 do not transfer
\item Within this scope the estimate is locally unique (this work; counterpart of Gy\"or\"ok 2026 Sec.~4.2)
\item Unique is not true: error $\Phi_{\mathcal R}^{+}\delta^\star$, recovery if and only if $\Phi_{\mathcal R}^\top\delta^\star=0$ (this work; counterpart of Gy\"or\"ok 2026 Cond.~4, Thm.~7, Eq.~(21))
\item Unique beats non-unique in the worst case; the display is a one-step surrogate of the trained model (Gy\"or\"ok 2026 Sec.~4.2)
\item Whether the condition holds depends on omitted dynamics and excitation; it needs $\theta^\star_{\mathrm{base}}$, so it is checked in simulation only, by the overlap $\gamma$ (Gy\"or\"ok 2026 Remark~5; this work)
\end{itemize}
\fi

The orthogonality of \eqref{eq:obc_coef} is exact but narrow.
It holds for the stack over $\mathcal R$, not at each tuple.
It holds for one step of the transition, at the expansion point of the current
epoch.
Its inner product is the Euclidean one in the normalised coordinates of
\eqref{eq:aug_norm}, so the normalisation sets its metric.
The tuples lie on $\bar x=0$, so the part of $f_{\mathrm{aug}}$ that acts
through nonzero added states is not protected.
After PS2 the added states are nonzero along the trajectories, so this route
is open in $\mathcal M_{\mathrm{OBC}}$.
Nothing constrains the multistep or closed-loop response.
The consistency and zero-covariance results
of~\cite[Thms.~16, 18]{gyorok2026obc}, derived for one-step prediction with
input-output models linear in their parameters, therefore do not transfer.

Within this scope, we show that the estimate is unique.

\noindent\textit{Proposition 1 (local uniqueness):} Let
$\Phi_{\mathcal R}$ have full column rank, and let $\tilde F$ and
$\tilde F'$ be corrected stacks
$F_{\mathcal R}-\Phi_{\mathcal R}\theta_{\mathrm{aux}}$ of
\eqref{eq:obc_coef}. If the linearised stacked transitions of
$(\bar\theta_{\mathrm{base}}+p,\tilde F')$ and
$(\bar\theta_{\mathrm{base}},\tilde F)$ on $\mathcal R$ are equal, then $p=0$.
\begin{IEEEproof}
Equality gives $\Phi_{\mathcal R}p=\tilde F-\tilde F'$. The right-hand side is
orthogonal to $\ran\Phi_{\mathcal R}$ and the left-hand side lies in it, so
$\Phi_{\mathcal R}p=0$, and full column rank gives $p=0$.
\end{IEEEproof}
Without the correction, every $p$ is matched by the network term
$-\Phi_{\mathcal R}p$ of \eqref{eq:obc_negation}.
The statement is local and concerns the reference stack.
It does not show that $V$ has a unique minimiser in $\theta_{\mathrm{base}}$.

A unique estimate need not be the true one.
Let $\theta^\star_{\mathrm{base}}$ be the true combinations and
\begin{equation}
 \delta^\star=\col\bigl(\tilde x^{\mathrm r}_{k+1}
 -f_{\mathrm{base}}(\theta^\star_{\mathrm{base}},\tilde x^{\mathrm r}_k,u_k)\bigr)_{k\in\mathcal K}
 \in\R^{6N_{\mathcal R}}
\label{eq:obc_delta}
\end{equation}
the true one-step discrepancy on $\mathcal R$, where the successor
$\tilde x^{\mathrm r}_{k+1}$ is reconstructed with the same stencil.
Suppose the data are noise-free and the corrected model reproduces every
successor on $\mathcal R$.
Suppose further that $f_{\mathrm{base}}$ is affine in $\theta_{\mathrm{base}}$
on the segment from $\theta^\star_{\mathrm{base}}$ to the estimate
$\hat\theta_{\mathrm{base}}$, so that $\Phi_{\mathcal R}$ is the same along it.
Then $\Phi_{\mathcal R}(\hat\theta_{\mathrm{base}}-\theta^\star_{\mathrm{base}})
+\tilde F=\delta^\star$, and multiplying by $\Phi_{\mathcal R}^{+}$ removes
$\tilde F$:
\begin{equation}
 \hat\theta_{\mathrm{base}}-\theta^\star_{\mathrm{base}}
 =\Phi_{\mathcal R}^{+}\delta^\star,
 \qquad
 \hat\theta_{\mathrm{base}}=\theta^\star_{\mathrm{base}}
 \;\Longleftrightarrow\;
 \Phi_{\mathcal R}^{\top}\delta^\star=0.
\label{eq:obc_bias}
\end{equation}
This is the state-level counterpart of~\cite[Cond.~4, Thm.~7,
Eq.~(21)]{gyorok2026obc}.

Where the condition fails, the estimate stays unique with error
$\Phi_{\mathcal R}^{+}\delta^\star$, whereas the additive structure without
correction can give larger errors in the worst
case~\cite[Sec.~4.2]{gyorok2026obc}.
The training criterion is a closed-loop multistep output error, so
\eqref{eq:obc_bias} is a one-step surrogate, not the bias of the trained
model.

Whether the condition holds depends on the omitted dynamics and the
excitation.
When the structure of the omitted term is known, the input can be designed to
meet it~\cite[Remark~5]{gyorok2026obc}.
The records of Section~\ref{sec:setup} are not designed for it.
The condition also involves $\theta^\star_{\mathrm{base}}$, so it can be
checked in simulation only, by the overlap
\begin{equation}
 \gamma(v;\Phi_{\mathcal S})
 =\frac{\norm{\Phi_{\mathcal S}\Phi_{\mathcal S}^{+}v}}{\norm{v}}\in[0,1],
\label{eq:obc_overlap}
\end{equation}
where $\Phi_{\mathcal S}\in\R^{m\times10}$ is a stack of sensitivities over a
tuple set $\mathcal S$ and $v\in\R^m$ a stacked vector on the same tuples.
$\gamma(\delta^\star;\Phi_{\mathcal R})=0$ exactly when the condition holds,
and for the corrected network on held-out tuples, $\gamma$ shows how far the
orthogonality extends beyond $\mathcal R$.
\todo{Missing: where the thesis states that the simulated truth violates the condition. Candidate: Section~VII, since a friction force $-c\,\mathrm{sign}(\dot q)$ has the nonzero inner product $-c\sum_k|\dot q_k|$ with the damping column of the same axis on every moving record (docs/references.md, gyorok2026obc row, Condition~4 analysis), so $\Phi_{\mathcal R}^\top\delta^\star\neq0$.}
\todo{Missing: notation in later sections. Section~VI uses $\rho^*$ and $\rho_f$, Section~VII and the appendix use $J$ and $\Delta^*$; this section uses $\Phi_{\mathcal R}$, $\delta^\star$ and $\gamma$, because $\rho$ is the margin of \eqref{eq:mdelta_map} and $\Delta$ the scheduling block. Candidate: adopt this section's symbols there. THESIS-RESULTS step~5 lists only combination errors, so whether Results reports $\gamma(\delta^\star;\Phi_{\mathcal R})$ is Dirk's decision.}
